\documentclass[a4paper,12pt]{extarticle}
\usepackage{geometry}
\usepackage[T1]{fontenc}
\usepackage[utf8]{inputenc}
\usepackage[english,russian]{babel}
\usepackage{amsmath}
\usepackage{amsthm}
\usepackage{amssymb}
\usepackage{fancyhdr}
\usepackage{setspace}
\usepackage{graphicx}
\usepackage{colortbl}
\usepackage{tikz}
\usepackage{pgf}
\usepackage{subcaption}
\usepackage{listings}
\usepackage{indentfirst}
\usepackage[
backend=biber,
style=numeric,
maxbibnames=99
]{biblatex}
\addbibresource{refs.bib}
\usepackage[colorlinks,citecolor=blue,linkcolor=blue,bookmarks=false,hypertexnames=true, urlcolor=blue]{hyperref} 
\usepackage{indentfirst}
\usepackage{mathtools}
\usepackage{booktabs}
\usepackage[flushleft]{threeparttable}
\usepackage{tablefootnote}

\usepackage{chngcntr} % нумерация графиков и таблиц по секциям
\counterwithin{table}{section}
\counterwithin{figure}{section}

\graphicspath{{graphics/}}%путь к рисункам

\makeatletter
% \renewcommand{\@biblabel}[1]{#1.} % Заменяем библиографию с квадратных скобок на точку:
\makeatother

\geometry{left=2.5cm}% левое поле
\geometry{right=1.0cm}% правое поле
\geometry{top=2.0cm}% верхнее поле
\geometry{bottom=2.0cm}% нижнее поле
\setlength{\parindent}{1.25cm}
\renewcommand{\baselinestretch}{1.5} % междустрочный интервал


\newcommand{\bibref}[3]{\hyperlink{#1}{#2 (#3)}} % biblabel, authors, year
\addto\captionsrussian{\def\refname{Список литературы (или источников)}} 

\renewcommand{\theenumi}{\arabic{enumi}}% Меняем везде перечисления на цифра.цифра
\renewcommand{\labelenumi}{\arabic{enumi}}% Меняем везде перечисления на цифра.цифра
\renewcommand{\theenumii}{.\arabic{enumii}}% Меняем везде перечисления на цифра.цифра
\renewcommand{\labelenumii}{\arabic{enumi}.\arabic{enumii}.}% Меняем везде перечисления на цифра.цифра
\renewcommand{\theenumiii}{.\arabic{enumiii}}% Меняем везде перечисления на цифра.цифра
\renewcommand{\labelenumiii}{\arabic{enumi}.\arabic{enumii}.\arabic{enumiii}.}% Меняем везде перечисления на цифра.цифра
\begin{document}

\maketitle

\section{Представление музыкальных данных в компьютерных системах}
\label{sec:music_representation}

Представление музыки в компьютерных системах является одним из ключевых вопросов для задач анализа, синтеза и генерации музыкального материала, поскольку именно выбранная форма данных определяет, какие свойства музыки будут доступны модели напрямую, а какие --- только косвенно \cite{briot2017deep,ji2020comprehensive,duan2025recent}. 
В современной литературе по генерации музыки музыкальные данные обычно рассматриваются как многоуровневые и мультимодальные, а различия между форматами представления связываются как с музыкальной природой данных, так и с вычислительными ограничениями конкретных методов \cite{ji2020comprehensive,duan2025recent}. 

\subsection{Основные типы представления музыки}
\label{subsec:main_types_repr}

В большинстве современных работ выделяются два основных типа представления музыкальных данных: аудиальное и символьное \cite{briot2017deep,ji2020comprehensive,duan2025recent}. 
Аудиальное представление описывает музыку как звуковой сигнал, то есть как непрерывную акустическую реализацию произведения \cite{briot2017deep,ji2020comprehensive}. 
Символьное представление, напротив, описывает музыку как набор дискретных музыкальных событий или абстрактных структур, например нот, длительностей, тактовых позиций, темпа и аккордов \cite{ji2023symbolic,duan2025recent,ji2020comprehensive}. 
Такое разделение важно для задач генерации музыки, поскольку аудиальные данные лучше сохраняют тембр и нюансы звучания, а символьные --- композиционную структуру произведения \cite{ji2020comprehensive,duan2025recent,ji2023symbolic}. 

\subsection{Аудиальное представление музыки}
\label{subsec:audio_representation}

Аудиальное представление музыки является непрерывным и ориентировано прежде всего на акустические свойства сигнала, включая тембр, спектральную структуру и микродинамику \cite{ji2020comprehensive,duan2025recent}. 
В обзорах по deep music generation аудио-представления обычно делятся на одномерные и двумерные \cite{ji2020comprehensive}. 
К первым относится временная волновая форма сигнала, ко вторым --- различные спектральные представления \cite{briot2017deep,ji2020comprehensive}. 

\subsubsection{Волновая форма сигнала}
\label{subsubsec:waveform}

Наиболее прямой формой представления музыки в компьютере является waveform, то есть последовательность дискретных отсчётов звукового сигнала во времени \cite{briot2017deep,ji2020comprehensive,duan2025recent}. 
Такое представление является наиболее близким к физической природе звука и сохраняет максимум акустической информации, содержащейся в записи \cite{ji2020comprehensive,duan2025recent}. 
Именно поэтому ряд моделей аудиогенерации использует исходную волновую форму в качестве входных данных \cite{ji2020comprehensive}. 
В то же время работа с waveform связана с высокой вычислительной стоимостью, поскольку даже короткий музыкальный фрагмент содержит очень большое число отсчётов \cite{briot2017deep,ji2020comprehensive}. 
Например, при частоте дискретизации 44{,}1 кГц стереофонический сигнал длительностью одна минута содержит
\[
44100 \times 60 \times 2 = 5\,292\,000
\]
отсчётов, что иллюстрирует высокую размерность такого представления. 

\subsubsection{Спектрограмма}
\label{subsubsec:spectrogram}

Для уменьшения размерности и выделения более устойчивых признаков музыкального сигнала аудио часто переводят из временной области в частотную или временно-частотную \cite{briot2017deep,ji2020comprehensive,duan2025recent}. 
Наиболее распространённым представлением такого рода является спектрограмма, обычно получаемая с помощью кратковременного преобразования Фурье \cite{briot2017deep,ji2020comprehensive,duan2025recent}. 
Спектрограмма представляет собой временно-частотную матрицу, в которой одна ось соответствует времени, другая --- частоте, а значение ячейки отражает энергию сигнала на данной частоте в данный момент \cite{duan2025recent,ji2020comprehensive}. 
По сравнению с исходной волновой формой спектрограмма удобнее для обработки нейросетями, особенно сверточными архитектурами \cite{ji2020comprehensive}. 
При этом в амплитудной спектрограмме фазовая информация выражена неполно, поэтому восстановление сигнала требует дополнительных процедур фазовой реконструкции \cite{ji2020comprehensive}. 

\subsubsection{Mel-спектрограмма}
\label{subsubsec:mel_spec}

Одним из наиболее используемых вариантов спектрального представления является mel-спектрограмма \cite{duan2025recent,ji2020comprehensive}. 
В этом случае частотная ось обычной спектрограммы проецируется на шкалу Мела, которая приближает восприятие частоты человеческим слухом \cite{duan2025recent,ji2020comprehensive}. 
Такое преобразование делает представление более перцептивно согласованным, поскольку слух человека различает частоты нелинейно \cite{duan2025recent}. 
Именно поэтому mel-спектрограммы широко используются в задачах генерации и анализа музыкального аудио \cite{duan2025recent,ji2020comprehensive}. 

\subsubsection{Хромаграмма и родственные представления}
\label{subsubsec:chromagram}

Для музыкальных задач, где особенно важны гармония, тональность и распределение высотных классов, применяется хромаграмма \cite{duan2025recent,ji2020comprehensive}. 
Хромаграмма агрегирует частотную информацию по двенадцати полутонам хроматической шкалы и тем самым описывает прежде всего гармоническое содержание сигнала \cite{duan2025recent}. 
Такое представление полезно в задачах моделирования аккордов, тонального контекста и структуры произведения \cite{duan2025recent,ji2020comprehensive}. 

\subsection{Символьное представление музыки}
\label{subsec:symbolic_representation}

Символьное представление музыки описывает не сам акустический сигнал, а музыкальную структуру в виде дискретных сущностей \cite{ji2023symbolic,duan2025recent,ji2020comprehensive}. 
В этом случае музыка задаётся через ноты, длительности, темп, позиции в такте, аккорды, инструменты и другие параметры, существенные для композиционной организации \cite{ji2023symbolic,duan2025recent}. 
Такой подход ближе к нотной записи и особенно удобен в задачах генерации мелодии, гармонизации, аранжировки и многодорожечной композиции \cite{ji2023symbolic,ji2020comprehensive}. 

\subsubsection{Piano roll}
\label{subsubsec:pianoroll}

Piano roll представляет музыку в виде двумерной матрицы, где вертикальная ось соответствует высотам нот, горизонтальная --- времени, а значения ячеек показывают наличие ноты и, в ряде вариантов, её динамику \cite{duan2025recent,ji2023symbolic,briot2017deep}. 
Такой формат нагляден, явно отражает метрическую сетку и хорошо согласуется с интуитивным визуальным пониманием музыкальной структуры \cite{duan2025recent,ji2023symbolic}. 
Благодаря матричной форме piano roll часто оказывается удобным для моделей, работающих с двумерными структурами данных \cite{briot2017deep,ji2023symbolic}. 
Однако стандартный piano roll не содержит явной информации о моменте окончания ноты, из-за чего возникает неоднозначность между одной длинной нотой и несколькими повторяющимися одинаковыми нотами \cite{duan2025recent}. 

\subsubsection{ABC-нотация}
\label{subsubsec:abc}

ABC-нотация является текстовым форматом музыкальной записи, в котором высоты кодируются буквами, а ритмические и структурные характеристики --- дополнительными символами \cite{duan2025recent,ji2023symbolic}. 
Такое представление компактно и удобно для обработки методами последовательностного моделирования, поскольку музыка в нём становится похожей на текст \cite{briot2017deep,ji2023symbolic}. 
ABC-нотация особенно удобна для относительно простых мелодических корпусов и традиционной монодической музыки \cite{duan2025recent,ji2023symbolic}. 
При этом по сравнению с MIDI-подобными форматами она хуже подходит для сложной полифонии, многодорожечной фактуры и детального описания исполнительских характеристик \cite{duan2025recent,ji2023symbolic}. 

\subsubsection{MIDI как базовое символьное представление}
\label{subsubsec:midi}

MIDI (Musical Instrument Digital Interface) является наиболее распространённым символьным форматом в задачах компьютерной музыки \cite{briot2017deep,duan2025recent,ji2020comprehensive}. 
MIDI представляет музыку не в виде звуковой волны, а как последовательность событий исполнения, например \texttt{Note-On}, \texttt{Note-Off}, изменения управляющих параметров и временные сдвиги между событиями \cite{briot2017deep,duan2025recent}. 
Иными словами, MIDI хранит инструкции о том, какие ноты должны быть сыграны, когда именно они начинаются, как долго звучат и с какой динамикой исполняются \cite{duan2025recent,briot2017deep}. 
Такой подход делает MIDI-файлы существенно компактнее аудиофайлов, поскольку они содержат не сами звуковые отсчёты, а только описание музыкальных действий \cite{duan2025recent}. 
Кроме того, MIDI хорошо совместим с электронными инструментами и программными средствами, а одна и та же MIDI-последовательность может быть озвучена разными тембрами \cite{duan2025recent,briot2017deep}. 

\subsubsection{Улучшения и расширения MIDI-представления}
\label{subsubsec:midi_extensions}

Поскольку стандартного MIDI-представления часто оказывается недостаточно для современных генеративных моделей, в литературе были предложены различные расширенные токенизации MIDI-подобных событий \cite{duan2025recent,ji2023symbolic}. 
Общее направление этих работ состоит в том, чтобы сохранить компактность event-based представления, но сделать более явными метрическую, ритмическую, гармоническую и инструментальную структуру музыки \cite{ji2023symbolic,duan2025recent}. 
Одним из наиболее известных расширений является REMI (REvamped MIDI-derived events) \cite{ji2023symbolic,duan2025recent}. 
В REMI стандартный поток MIDI-событий дополняется специальными токенами, кодирующими границы тактов, позиции внутри такта, темп и аккордовую информацию \cite{ji2023symbolic,duan2025recent}. 
Это позволяет сделать метрическую сетку явной и облегчает моделирование ритма и формы произведения \cite{ji2023symbolic}. 
Дальнейшим развитием этого подхода стал формат REMI+, который вводит дополнительные токены размера и инструмента \cite{duan2025recent}. 
Помимо REMI и REMI+, в литературе предлагаются и другие способы улучшения MIDI-подобных представлений, включая MuMIDI, Compound Word и OctupleMIDI \cite{ji2023symbolic,duan2025recent}. 

\subsection{Выводы по представлению музыкальных данных}
\label{subsec:repr_conclusion}

Таким образом, в компьютерных системах музыка может быть представлена либо как аудиосигнал, либо как символьная последовательность \cite{briot2017deep,duan2025recent,ji2020comprehensive}. 
Аудиальные представления лучше сохраняют тембровую и акустическую полноту сигнала, но требуют больших вычислительных ресурсов \cite{ji2020comprehensive,duan2025recent}. 
Символьные представления компактнее, структурнее и лучше подходят для моделирования композиционной логики произведения \cite{ji2023symbolic,duan2025recent}. 
Поэтому выбор формы представления данных является одним из фундаментальных факторов, влияющих на архитектуру модели, способ обучения и качество результата \cite{briot2017deep,ji2020comprehensive,duan2025recent}. 
\end{document}